% !TEX program = xelatex
% !BIB TS-program = biber

% **************************************************
% Document Class Definition
% **************************************************
\documentclass[%
    paper=A4,               % paper size --> A4 is default in Germany
    twoside=true,           % onesite or twoside printing
    openany,              % doublepage cleaning ends up right side
    parskip=full,           % spacing value / method for paragraphs
    chapterprefix=true,     % prefix for chapter marks
    11pt,                   % font size
    headings=normal,        % size of headings
    bibliography=totoc,     % include bib in toc
    listof=totoc,           % include listof entries in toc
    titlepage=on,           % own page for each title page
    captions=tableabove,    % display table captions above the float env
    draft=false,            % value for draft version
]{scrreprt}

\newcommand{\reportTitle}{\textcolor{ugent_blue}{Constrained Unrestricted Hartree Fock and Single Excitation Configuration Interaction}} 
\newcommand{\reportName}{\textbf{Ruben Van der Stichelen}}
\newcommand{\reportSubject}{\emph{Programming project}}
\newcommand{\reportDate}{}
\newcommand{\reportVersion}{}

\usepackage[utf8]{inputenc}
\usepackage[english]{babel} % babel system, adjust the language of the content
\PassOptionsToPackage{% setup clean thesis style
    figuresep=colon,%
    sansserif=true,%
    hangfigurecaption=false,%
    hangsection=true,%
    hangsubsection=true,%
    colorize=full,%
    colortheme=bluemagenta,%
    bibsys=biber,%
    bibfile=report,%
    bibstyle=numeric-comp,%
    wrapfooter=false,%
}{cleanthesis}
\usepackage{cleanthesis}

\hypersetup{% setup the hyperref-package options
    pdftitle={\reportTitle},    %   - title (PDF meta)
    pdfsubject={\reportSubject},%   - subject (PDF meta)
    pdfauthor={\reportName},    %   - author (PDF meta)
    plainpages=false,           %   -
    colorlinks=false,           %   - colorize links?
    pdfborder={0 0 0},          %   -
    breaklinks=true,            %   - allow line break inside links
    bookmarksnumbered=true,     %
    bookmarksopen=true          %
}

\definecolor{ugent_blue}{RGB}{30, 100, 200}
\definecolor{sciences}{RGB}{45, 140, 168}

\usepackage{algorithm}              % http://ctan.org/pkg/algorithms
\usepackage{algpseudocode}          % http://ctan.org/pkg/algorithmicx
\usepackage{amsmath}
\usepackage{amsfonts}             % mathfrak
\usepackage{amssymb}              % \mathbb
\usepackage[mathscr]{euscript}    % \mathscr
\usepackage{chemmacros}
\usepackage{float}
\usepackage{graphicx}
\usepackage{interval}
\usepackage{mathtools}            % dcases, DeclareMathOperator, DeclairePairedDelimiter
\usepackage{physics}               % physics typesetting, especially bra-ket notation
\usepackage{siunitx}
\usepackage{subcaption}
\usepackage{tensor}
\usepackage{xcolor}
\usepackage{ifxetex}
\usepackage{multicol}
\usepackage{pdfpages}
\usepackage{epigraph}
\usepackage{bm}
\usepackage[normalem]{ulem}
\usepackage[version=3]{mhchem}
\usepackage[linewidth=1pt]{mdframed}
\usepackage{enumitem}


% Default fixed font does not support bold face
\DeclareFixedFont{\ttb}{T1}{txtt}{bx}{n}{12} % for bold
\DeclareFixedFont{\ttm}{T1}{txtt}{m}{n}{12}  % for normal

% Custom colors
\usepackage{color}
\definecolor{deepblue}{rgb}{0,0,0.5}
\definecolor{deepred}{rgb}{0.6,0,0}
\definecolor{deepgreen}{rgb}{0,0.5,0}
\usepackage{listings}

% Python style for highlighting
\newcommand\pythonstyle{\lstset{
language=Python,
basicstyle=\ttm,
morekeywords={self},              % Add keywords here
keywordstyle=\ttb\color{deepblue},
emph={MyClass,__init__},          % Custom highlighting
emphstyle=\ttb\color{deepred},    % Custom highlighting style
stringstyle=\color{deepgreen},
frame=tb,                         % Any extra options here
showstringspaces=false
}}


% Python environment
\lstnewenvironment{python}[1][]
{
\pythonstyle
\lstset{#1}
}
{}

% Python for external files
\newcommand\pythonexternal[2][]{{
\pythonstyle
\lstinputlisting[#1]{#2}}}

% Python for inline
\newcommand\pythoninline[1]{{\pythonstyle\lstinline!#1!}}

\DeclareSymbolFont{euleroperators}{U}{eur}{m}{n}
\SetSymbolFont{euleroperators}{bold}{U}{eur}{b}{n}

\makeatletter
\renewcommand{\operator@font}{\mathgroup\symeuleroperators}
\makeatother

% \usepackage[utf8]{inputenc} 


\ifxetex
  \usepackage{fontspec}
	% \setmainfont[
	% Path = ./,
	% BoldFont = UGentPannoText-SemiBold.otf,
	% ItalicFont = UGentPannoText-SemiLight.otf,
	% BoldItalicFont  = UGentPannoText-Medium.otf]
	% {UGentPannoText-Normal.otf}
	\newfontfamily\PannoSemiLight{UGentPannoText-SemiLight.otf}
	\newfontfamily\PannoMedium{UGentPannoText-Medium.otf}
	\newfontfamily\PannoSemiBold{UGentPannoText-SemiBold.otf}
	\newcommand{\PannoLarge}{\LARGE}
    \newcommand{\PannoNormalsize}{\large}
    \renewcommand{\helv}{\fontspec{UGentPannoText-Normal.otf}\fontsize{9}{11}\selectfont}
    \renewcommand{\book}{\fontspec{UGentPannoText-Normal.otf}\fontsize{11}{13}\selectfont}
    \renewcommand{\tgherosfont}{\fontspec{UGentPannoText-Normal.otf}}
    \renewcommand{\thesischapterfont}{\color{ctcolorblack}\huge\fontspec{UGentPannoText-Normal.otf}}
\else
  % \usepackage[utf8]{inputenc}           % Om niet ascii karakters rechtstreeks te kunnen typen
	\usepackage{lmodern}
	\usepackage[T1]{fontenc}
	\renewcommand{\familydefault}{\sfdefault}
	\newcommand{\PannoSemiLight}{}
	\newcommand{\PannoMedium}{}
	\newcommand{\PannoSemiBold}{}
	\newcommand{\PannoLarge}{\Large}
	\newcommand{\PannoNormalsize}{\normalsize}
\fi

\graphicspath{{img/}}
\linespread{1.2}
\usepackage{chngcntr}
\counterwithout{section}{chapter}
\counterwithout{equation}{chapter}
\counterwithout{figure}{chapter}
\counterwithout{table}{chapter}


\title{\reportTitle}
\author{\reportName\\GQCG}

% **************************************************
% Document CONTENT
% **************************************************
\begin{document}

    \begin{titlepage}
        \hspace{-1cm}\includegraphics[height=3cm]{faculty.png}
        \begin{center}
            \vspace{1cm}
            \Huge{\reportTitle} \\
            \vspace{1cm}
            \Large{\reportName} \\
            \vspace{2cm}
            \includegraphics[height=2cm]{GQCG.png} \\
            \vspace{0.5cm}
            \large{Ghent University} \\ 
            \large{\today}
        \end{center}
        \vspace{4cm}
        \hspace{-1cm}\includegraphics[height=4cm]{UGent.png}
    \end{titlepage}
    % --------------------------
    % rename document parts
    % --------------------------
    \renewcaptionname{english}{\figurename}{Figure}
    \renewcaptionname{english}{\tablename}{Table}

    \setcounter{tocdepth}{2}		% define depth of toc

    % --------------------------
    % Body matter / main matter
    % --------------------------
    \pagenumbering{arabic}			% arabic page numbering
    \setcounter{page}{1}			% set page counter
    \pagestyle{maincontentstyle} 	% fancy header and footer

    % The following input statements correspond to the different chapters of these notes
    \section{Introduction}
\label{sec:intro}
In this paper, we will discuss the effect of CIS (configuration interaction singles) on the CUHF (constrained unrestricted hartree fock theory) wave function. We will do this in three
steps. First we will discuss the CUHF theory and implementation, using python. Then we will follow the same approach for the theory and code behind CIS. Finally we will discuss what
we found, and talk about some possible application fields.


\section{Constrained Unrestricted Hartree Fock Theory}
\label{sec:cuhf}
\subsection{Theory}
\label{subsec:cuhf_theory}
\subsection{Implementation}
\label{subsec:cuhf_impl}
\subsection{Results}
\label{subsec:cuhf_res}

\section{Single Excitation Configuration Interaction}
\subsection{Theory}
\label{subsec:cis_theory}
\subsection{Implementation}
\label{subsec:cis_impl}

\section{The CUHF Wave Function in the CIS Approach}
\label{sec:combination}

\section{Applications}
\label{sec:applic}
    \clearpage

    % --------------------------
    % Back matter
    % --------------------------

    {%
    \setstretch{1.1}
    \renewcommand{\bibfont}{\normalfont\small}
    \setlength{\biblabelsep}{0pt}
    \setlength{\bibitemsep}{0.5\baselineskip plus 0.5\baselineskip}
    \printbibliography[nottype=online]
    \printbibliography[heading=subbibliography,title={Webpages},type=online,prefixnumbers={@}]
    }
    
    % **************************************************
    % End of Document CONTENT
    % **************************************************
    \end{document}
